% Libro 3 · Mundo 1 · Estrategias de lectura (app: pu-estrategias)
\mundo{Estrategias de lectura}{pu-estrategias}{Antes de responder hay que leer, y antes de leer conviene saber a qué se
va. Este mundo enseña a prepararse para un texto en treinta segundos, a construir su mapa mientras se lee y a administrar
el tiempo del examen, incluidas esas preguntas que piden la opción que \emph{no} corresponde.}{%
\item \textbf{1.1} Prelectura: mirar antes de leer
\item \textbf{1.2} Idea central y mapa del texto
\item \textbf{1.3} Tiempo y preguntas EXCEPTO}

% ------------------------------------------------------------------
\leccion{pu-el-1}{Prelectura}{Juego de subrayar en el texto lo que responde.}

\begin{reto}
Lee solo este fragmento: «Hay oficios que no mueren de golpe, sino por goteo: un cliente menos cada mes, una pieza que ya
nadie fabrica, un aprendiz que nunca llega». ¿De qué tratará, muy probablemente, el texto completo?
\opcionesl{De la historia de la industria en Colombia.}{De la desaparición lenta de un oficio tradicional.}{De cómo conseguir aprendices para un negocio.}{De la crisis económica de un país.}
\end{reto}

\tecnica{La técnica: treinta segundos que ahorran cinco minutos}
Leer un texto «a ciegas», de la primera a la última palabra, obliga a la memoria a guardar todo por si acaso. La
\clave{prelectura} da un mapa antes de entrar: con el mapa, la lectura se vuelve una comprobación y no una adivinanza.

Se miran, en este orden:
\begin{enumerate}
  \item El \textbf{título} y el subtítulo: anuncian el tema y a veces la postura.
  \item La \textbf{fuente} o el tipo de texto (si aparece): una crónica, un ensayo, una carta.
  \item La \textbf{primera frase de cada párrafo}: en los textos expositivos y argumentativos suele contener la idea del
  párrafo.
  \item El \textbf{último párrafo}: allí suele estar la conclusión o la tesis.
\end{enumerate}

Con eso se formulan dos o tres preguntas propias («¿qué le pasó al relojero?», «¿por qué se acaba el oficio?») que la
lectura completa va a responder.

\begin{regla}
\begin{itemize}
  \item Título $\to$ fuente $\to$ primeras frases $\to$ último párrafo. Treinta segundos.
  \item Hazle dos preguntas al texto antes de leerlo.
  \item Si el examen te permite, lee primero los enunciados de las preguntas (no las opciones): te dicen qué buscar.
\end{itemize}
\end{regla}

\begin{texto}{El último relojero de la calle Junín}{Crónica · texto original escrito para este libro}
\pa Hay oficios que no mueren de golpe, sino por goteo: un cliente menos cada mes, una pieza que ya nadie fabrica, un
aprendiz que nunca llega. Don Aurelio Montoya lleva cincuenta y dos años arreglando relojes en un local de dos metros de
ancho, en el centro de Medellín, y ha visto cómo el goteo se volvía silencio.

\pa Cuando abrió, a comienzos de los años setenta, en la misma cuadra había otros cinco relojeros. Los sábados, dice, la
fila llegaba hasta la esquina: obreros que traían el reloj de cuerda heredado del abuelo, señoras con relojes de pulsera
que se atrasaban, estudiantes que querían un cronómetro para las carreras. «Uno no daba abasto», recuerda, y señala una
caja de madera llena de manecillas sueltas, como si todavía esperaran a sus dueños.

\pa Los relojes de cuarzo, más baratos y casi infalibles, fueron el primer golpe. Ya no había que limpiar engranajes ni
ajustar resortes; bastaba cambiar una pila. Luego llegaron los teléfonos celulares, y con ellos una pregunta que don
Aurelio escucha con frecuencia: «¿Y para qué un reloj, si la hora está en el celular?». Él responde siempre lo mismo: que
el celular dice la hora, pero el reloj la guarda.

\pa Hoy sus clientes son de dos clases. Unos son coleccionistas que pagan bien por restaurar una pieza antigua y no tienen
prisa. Otros son hijos o nietos que llegan con el reloj de un difunto y una petición menos técnica que afectiva: «Que
vuelva a andar». A esos, confiesa, a veces no les cobra.

\pa Lo que más le preocupa no es la falta de clientes, sino la falta de manos. Ha tenido tres aprendices; ninguno se quedó.
El oficio exige años de paciencia antes de ganar bien, y los muchachos, dice sin reproche, «tienen afán». Cuando él se
retire, el local probablemente se convertirá en otra venta de accesorios para celular.

\pa Mientras tanto, cada mañana a las ocho levanta la reja, enciende la lámpara de cuello largo y se pone la lupa en el ojo
derecho. En la pared, treinta relojes marcan la misma hora con distintos sonidos. Ninguno se atrasa.
\end{texto}

\subsection{Preguntas comentadas}
\begin{comentada}{¿Cuál es la idea central del texto?}
\mal{A}{Los relojes de cuarzo son mejores que los de cuerda.}{\emph{No lo dice}: el texto menciona que son «más baratos y
casi infalibles» \pf{3}, pero no los compara como mejores ni es su tema.}
\bien{B}{Un relojero tradicional sigue trabajando mientras su oficio desaparece poco a poco.}{Recoge el «goteo» del \S\,1,
las causas \pf{3, 5} y el final \pf{6}. Es la idea que conecta todos los párrafos.}
\mal{C}{Los jóvenes no quieren trabajar.}{\emph{Dice de más}: el texto habla de tres aprendices que «tienen afán»
\pf{5}, y aclara que lo dice «sin reproche».}
\mal{D}{Los celulares acabaron con todos los oficios del centro de Medellín.}{\emph{Dice de más}: generaliza a «todos los
oficios» algo que el texto dice de uno.}
\end{comentada}

\begin{comentada}{Según el texto, ¿qué piden los hijos o nietos que llevan el reloj de un difunto?}
\mal{A}{Que lo restauren para venderlo a un coleccionista.}{\emph{Dice otra cosa}: mezcla los dos tipos de clientes del
\S\,4.}
\mal{B}{Que no les cobren.}{\emph{Dice otra cosa}: no lo piden; es don Aurelio quien a veces decide no cobrarles.}
\bien{C}{Que el reloj vuelva a funcionar.}{Es la paráfrasis de «Que vuelva a andar» \pf{4}.}
\mal{D}{Que les enseñe el oficio.}{\emph{No lo dice}.}
\end{comentada}

\begin{trampa}
Escoger la opción que \emph{suena bien} o que repite palabras del texto sin comprobar que dice lo mismo. Las opciones B y A
de la segunda pregunta usan palabras del \S\,4, pero cambian quién hace qué.
\end{trampa}

\begin{atajo}
Para la idea central, la respuesta debe servir para \emph{todos} los párrafos. Si una opción solo resume uno, es un
detalle, no la idea central.
\end{atajo}

\taller
\piso{Calienta}{prelectura}
\begin{ejercicios}
\item Antes de leer el texto, ¿qué dos preguntas te sugería el título? Escríbelas y di si el texto las responde. \sol{Por ejemplo: ¿quién es el último relojero?, ¿por qué es el último? El texto las responde en los párrafos 1, 3 y 5.}
\item ¿Qué palabra al margen pondrías a cada párrafo? \sol{1 presentación (el goteo); 2 el pasado; 3 las causas; 4 los clientes de hoy; 5 la preocupación (no hay aprendices); 6 la rutina y el cierre.}
\end{ejercicios}
\piso{Como en el examen}{opción múltiple}
\begin{ejercicios}
\item En el texto, la frase «el celular dice la hora, pero el reloj la guarda» \pf{3} expresa que, para don Aurelio,
\opcionesl{el celular es más exacto que el reloj.}{el reloj tiene un valor que va más allá de dar la hora.}{los relojes guardan la hora aunque se acabe la pila.}{nadie debería usar celular.}
\sol{B. «Guardar» la hora sugiere memoria, permanencia, valor afectivo; el texto lo confirma con los relojes de los difuntos \pf{4}.}
\item ¿Cuál fue, según el texto, el primer golpe para el oficio?
\opcionesd{los celulares}{la falta de aprendices}{los relojes de cuarzo}{los coleccionistas}
\sol{C. \pf{3}: «fueron el primer golpe».}
\item El último párrafo cumple la función de
\opcionesl{dar una solución al problema del oficio.}{cerrar con una imagen de la rutina que persiste.}{criticar a los aprendices.}{presentar un nuevo problema.}
\sol{B. Describe la rutina diaria y termina con «Ninguno se atrasa», que contrasta con el oficio que se acaba.}
\end{ejercicios}

\begin{retoresuelto}
El fragmento habla de «oficios» que mueren «por goteo», con ejemplos de pérdida lenta: anuncia la desaparición lenta de
un oficio. Es exactamente el tema de la crónica. \textbf{Opción B.} La prelectura acertó antes de leer el resto.
\end{retoresuelto}
\bitacora

% ------------------------------------------------------------------
\leccion{pu-el-2}{Idea central y mapa del texto}{Juego del mapa: arrastrar cada párrafo a su función.}

\begin{reto}
Un texto tiene cuatro párrafos: el primero presenta una creencia común, el segundo muestra datos que la contradicen, el
tercero explica por qué la creencia es tan popular y el cuarto propone qué hacer. ¿Dónde está, probablemente, la postura
del autor?
\opcionesd{en el primero}{en el segundo y el cuarto}{solo en el tercero}{en ninguno}
\end{reto}

\tecnica{La técnica: una palabra por párrafo}
Un texto bien escrito es un edificio: cada párrafo cumple una \clave{función} (presentar, explicar, ejemplificar,
contrastar, objetar, concluir) y todos sostienen una \clave{idea central}. El \clave{mapa del texto} es la lista de esas
funciones, una palabra por párrafo, escrita al margen mientras se lee.

\figura{mapatexto}{El mapa del texto: una función por párrafo y una idea que las reúne.}

\subsubsection{Idea central, tema y título}
El \clave{tema} es de qué se habla («el olvido»); la \clave{idea central} es qué se dice de eso («olvidar es necesario para
recordar bien»). Una buena idea central es una frase completa, cubre todo el texto y no se queda en un detalle.

\begin{regla}
\begin{itemize}
  \item Al margen de cada párrafo: una palabra que diga su función.
  \item Tema $=$ de qué. Idea central $=$ qué se afirma de eso.
  \item La idea central debe poder «explicar» todos los párrafos; si deja uno por fuera, es incompleta.
  \item En los argumentativos, busca la tesis en el primer o el último párrafo.
\end{itemize}
\end{regla}

\begin{texto}{Elogio del olvido}{Ensayo de divulgación · texto original escrito para este libro}
\pa Solemos hablar del olvido como de una falla: una avería de la memoria que hay que combatir con listas, alarmas y
fotografías. Quien olvida una cita o un nombre se disculpa como si hubiera cometido una falta. Sin embargo, buena parte de
lo que la ciencia ha aprendido en el último siglo sugiere lo contrario: olvidar no es el enemigo de la memoria, sino una de
sus herramientas.

\pa Piénsese en lo que pasaría si recordáramos todo. Cada calle recorrida, cada conversación banal, cada número de
teléfono marcado una sola vez ocuparían el mismo lugar que el rostro de una madre o la fórmula aprendida para un examen.
Recordar sería como buscar una carta importante en una casa donde nunca se ha botado un papel. Los pocos casos estudiados
de personas con una memoria casi total describen esa experiencia no como un don, sino como una carga.

\pa El cerebro, en cambio, selecciona. Lo que se repite, lo que emociona, lo que se conecta con algo ya sabido, tiende a
quedarse; lo demás se desvanece. Desde finales del siglo XIX se sabe que lo recién aprendido se pierde muy rápido en las
primeras horas y luego más despacio, y que cada repaso hace la caída menos empinada. Durante el sueño, además, el cerebro
parece reforzar algunos recuerdos y debilitar otros, como un editor que decide qué páginas conservar.

\pa Nada de esto significa que todo olvido sea bueno. Hay enfermedades que borran lo que no debería borrarse, y hay
olvidos colectivos —de injusticias, de víctimas— que una sociedad no puede permitirse. Pero confundir esos olvidos con el
olvido cotidiano lleva a una idea equivocada: que la mejor memoria es la que guarda más.

\pa Para un estudiante, la conclusión es práctica. No se trata de leer diez veces el mismo capítulo la víspera del
examen, sino de repasar poco y a menudo, de relacionar lo nuevo con lo conocido y de dormir bien. La memoria no es un
depósito que se llena: es un jardín que se poda.
\end{texto}

\subsection{Preguntas comentadas}
\begin{comentada}{¿Cuál es la idea central del texto?}
\mal{A}{La memoria humana es muy limitada.}{\emph{Dice otra cosa}: el texto no habla de límites sino de selección.}
\bien{B}{Olvidar es, en buena medida, parte del buen funcionamiento de la memoria.}{Es la tesis del \S\,1 («una de sus
herramientas»), desarrollada en 2 y 3, matizada en 4 y aplicada en 5.}
\mal{C}{Todo olvido es positivo.}{\emph{Dice de más}, y el \S\,4 lo niega expresamente: «Nada de esto significa que todo
olvido sea bueno».}
\mal{D}{Los estudiantes deben dormir más.}{Es un detalle del \S\,5, no la idea que organiza el texto.}
\end{comentada}

\begin{comentada}{¿Qué función cumple el cuarto párrafo?}
\mal{A}{Presentar la tesis.}{La tesis está en el \S\,1.}
\bien{B}{Matizar la tesis para evitar una interpretación exagerada.}{Reconoce olvidos dañinos sin abandonar la idea
central: es una concesión.}
\mal{C}{Dar un ejemplo de la vida cotidiana.}{\emph{Dice otra cosa}: los ejemplos cotidianos están en el \S\,2.}
\mal{D}{Concluir con una recomendación.}{Eso lo hace el \S\,5.}
\end{comentada}

\begin{trampa}
Elegir como idea central la frase más llamativa o la conclusión práctica. «La memoria es un jardín que se poda» resume bien,
pero «los estudiantes deben dormir más» es solo una consecuencia.
\end{trampa}

\taller
\piso{Calienta}{el mapa}
\begin{ejercicios}
\item Escribe el mapa del texto: una palabra por párrafo. \sol{1 tesis; 2 ejemplo hipotético; 3 explicación científica; 4 concesión o matiz; 5 conclusión práctica.}
\item Diferencia el tema de la idea central. \sol{Tema: el olvido. Idea central: olvidar es una herramienta de la memoria, no solo una falla.}
\end{ejercicios}
\piso{Como en el examen}{opción múltiple}
\begin{ejercicios}
\item La comparación con «una casa donde nunca se ha botado un papel» \pf{2} sirve para mostrar que
\opcionesl{guardar todo dificultaría encontrar lo importante.}{las personas desordenadas olvidan más.}{la memoria funciona como un archivo de papel.}{hay que botar los papeles viejos.}
\sol{A.}
\item Según el \S\,3, ¿qué hace que un recuerdo tienda a quedarse?
\opcionesl{que sea reciente}{que se repita, emocione o se conecte con algo ya sabido}{que se escriba en una lista}{que se aprenda la víspera}
\sol{B.}
\item La frase final, «es un jardín que se poda», resume la idea de que la memoria
\opcionesl{crece sin control.}{necesita cuidados médicos.}{mejora cuando se elimina lo que sobra.}{se pierde con los años.}
\sol{C.}
\end{ejercicios}

\begin{retoresuelto}
El primero presenta una creencia común (no la del autor); el segundo la contradice con datos y el cuarto propone qué hacer:
ahí está la postura. El tercero explica la popularidad de la creencia, pero no es la postura. \textbf{Opción B.}
\end{retoresuelto}
\bitacora

% ------------------------------------------------------------------
\leccion{pu-el-3}{Tiempo y preguntas EXCEPTO}{Juego de tachar las opciones que el texto sí dice.}

\begin{reto}
«Según el texto, la biblioteca ofrece todo lo siguiente, EXCEPTO…». En esta pregunta, ¿qué opción hay que marcar?
\opcionesl{La que el texto menciona primero.}{La que el texto no menciona o contradice.}{La que parece más importante.}{La que se repite más veces.}
\end{reto}

\tecnica{La técnica: invertir la pregunta}
Las preguntas \clave{EXCEPTO} (o «¿cuál \emph{no}…?», «es falso que…») piden la única opción que el texto \emph{no}
sostiene. Son peligrosas porque el cerebro, por costumbre, busca la verdadera. El método es convertirlas en un ejercicio de
tachado: se verifica cada opción contra el texto, se tacha la que sí aparece, y la que queda es la respuesta.

\begin{center}\small
\begin{tabular}{@{}lll@{}}
\toprule
Opción & ¿Está en el texto? & \\\midrule
A & sí, \S\,2 & \sout{tachar}\\
B & sí, \S\,3 & \sout{tachar}\\
C & no aparece & \textbf{respuesta}\\
D & sí, \S\,4 & \sout{tachar}\\\bottomrule
\end{tabular}
\end{center}

\subsubsection{El tiempo}
En el examen hay aproximadamente un minuto por pregunta, después de leer el texto. Tres reglas:
\begin{itemize}
  \item Primera vuelta: responde lo que sabes en menos de un minuto; marca lo dudoso.
  \item No releas el texto entero para una pregunta: ve al párrafo indicado.
  \item Nunca dejes una pregunta en blanco si no hay penalización: descarta y elige.
\end{itemize}

\begin{regla}
\begin{itemize}
  \item EXCEPTO: tacha las que el texto sí dice; marca la que sobra.
  \item Escribe «NO» grande junto a las preguntas negativas para no olvidarlo a mitad de camino.
  \item Un minuto por pregunta; las dudosas, al final.
\end{itemize}
\end{regla}

\begin{texto}{Las bibliotecas que no cierran}{Artículo · texto original escrito para este libro}
\pa En muchos barrios de las ciudades colombianas, la biblioteca pública dejó hace tiempo de ser un lugar de silencio
absoluto. Quien entra un sábado a una de ellas puede encontrar, en la misma tarde, un taller de escritura para adultos
mayores, una clase de ajedrez para niños y un grupo de jóvenes preparando el examen de admisión a la universidad.

\pa Ese cambio no fue un accidente. Durante años, las bibliotecas vieron cómo bajaban los préstamos de libros mientras
crecía la demanda de otras cosas: un computador con internet, un lugar seguro para hacer tareas, un espacio para reunirse.
En lugar de resistirse, muchas decidieron ampliar su idea de lo que significa leer.

\pa Hoy ofrecen préstamo de libros y también de tabletas, salas de cómputo gratuitas, cursos para aprender a hacer trámites
en línea y asesoría para buscar empleo. Algunas tienen ludotecas, donde los más pequeños aprenden jugando, y otras han
abierto huertas comunitarias en sus terrazas.

\pa Los bibliotecarios cuentan que el cambio trajo tensiones. Hay usuarios que extrañan el silencio y otros que sienten que
los libros perdieron protagonismo. La respuesta más frecuente ha sido separar los espacios: salas silenciosas para quien
quiere leer o estudiar, y salas abiertas para las actividades ruidosas.

\pa Lo que no ha cambiado es la regla esencial: la entrada es libre y nadie tiene que demostrar nada para usar el lugar. En
barrios donde no hay parques ni centros culturales, esa puerta abierta convierte a la biblioteca en algo parecido a una
plaza con techo.
\end{texto}

\subsection{Preguntas comentadas}
\begin{comentada}{Según el texto, las bibliotecas ofrecen todo lo siguiente, EXCEPTO}
\mal{A}{préstamo de tabletas.}{Está en el \S\,3: se tacha.}
\mal{B}{asesoría para buscar empleo.}{Está en el \S\,3: se tacha.}
\bien{C}{atención médica gratuita.}{No aparece en ninguna parte del texto: es la respuesta.}
\mal{D}{clases de ajedrez para niños.}{Está en el \S\,1: se tacha.}
\end{comentada}

\begin{comentada}{¿Cuál de las siguientes afirmaciones NO se puede sostener a partir del texto?}
\mal{A}{Los préstamos de libros disminuyeron en algún momento.}{\S\,2: «vieron cómo bajaban los préstamos».}
\mal{B}{Algunos usuarios no están de acuerdo con los cambios.}{\S\,4: hay quienes extrañan el silencio.}
\mal{C}{Para entrar no hay que demostrar nada.}{\S\,5.}
\bien{D}{Las bibliotecas eliminaron las salas de silencio.}{Lo contrario: el \S\,4 dice que se separaron los espacios y que
hay salas silenciosas. \emph{Dice lo contrario}.}
\end{comentada}

\begin{trampa}
Olvidar a mitad de la pregunta que era negativa y marcar la primera opción verdadera que aparece. Por eso se escribe «NO» al
margen.
\end{trampa}

\taller
\piso{Calienta}{tachar}
\begin{ejercicios}
\item Haz la tabla de tachado para: «Según el \S\,3, las bibliotecas tienen todo lo siguiente EXCEPTO: A) huertas; B) ludotecas; C) salas de cine; D) salas de cómputo». \sol{A, B y D están en el \S\,3; C no aparece: respuesta C.}
\end{ejercicios}
\piso{Como en el examen}{opción múltiple}
\begin{ejercicios}
\item Según el texto, las bibliotecas cambiaron porque
\opcionesl{el gobierno lo ordenó.}{bajaron los préstamos y creció la demanda de otros servicios.}{los bibliotecarios querían más ruido.}{los libros se volvieron muy caros.}
\sol{B. \S\,2.}
\item La expresión «una plaza con techo» \pf{5} se refiere a que la biblioteca
\opcionesl{tiene una construcción antigua.}{funciona como un lugar público de encuentro para todos.}{se construyó en una plaza.}{no tiene libros.}
\sol{B.}
\item Todas las siguientes son tensiones o reacciones mencionadas en el \S\,4, EXCEPTO
\opcionesl{usuarios que extrañan el silencio.}{la sensación de que los libros perdieron protagonismo.}{la separación de espacios.}{el cierre de bibliotecas por falta de usuarios.}
\sol{D. No se menciona ningún cierre (el título dice lo contrario).}
\end{ejercicios}

\begin{retoresuelto}
En una pregunta EXCEPTO se marca la opción que el texto no sostiene. \textbf{Opción B.}
\end{retoresuelto}
\bitacora

% ------------------------------------------------------------------
\evaluacion
\begin{texto}{La ciudad que aprendió a mirar el cielo}{Crónica · texto original escrito para este libro}
\pa Durante décadas, en el pueblo de San Telmo nadie levantaba la vista de noche. No había razón: el alumbrado público,
instalado sin mucho cuidado, iluminaba más el cielo que las calles, y las estrellas se habían vuelto un rumor de los
abuelos.

\pa Todo cambió cuando una profesora de física del colegio, cansada de explicar las constelaciones con fotocopias, propuso
apagar las luces de la plaza durante una hora, una vez al mes. La alcaldía aceptó a regañadientes, con la condición de que
hubiera policías vigilando.

\pa La primera noche llegaron cuarenta personas. La tercera, más de trescientas. Hubo quienes vieron por primera vez la Vía
Láctea y quienes lloraron sin saber explicar por qué. Los niños aprendieron a encontrar a Orión; los viejos recordaron las
historias que les contaban sus padres sobre el cielo.

\pa El éxito trajo una discusión que nadie esperaba. Si una hora de oscuridad hacía tanto bien, ¿por qué no cambiar las
lámparas para que alumbraran hacia el suelo y no hacia arriba? Los comerciantes temían por la seguridad; los ingenieros
explicaron que unas lámparas bien dirigidas iluminan mejor la calle y gastan menos energía.

\pa Dos años después, San Telmo cambió la mitad de su alumbrado. Las noches de plaza oscura siguen, pero ya no son el único
momento en que se ven estrellas. La profesora dice que su mayor logro no fue enseñar astronomía, sino conseguir que un
pueblo entero se preguntara qué estaba perdiendo sin saberlo.
\end{texto}
\begin{ejercicios}
\item ¿Cuál es la idea central del texto?
\opcionesl{Las lámparas modernas gastan menos energía.}{Una iniciativa escolar llevó a un pueblo a recuperar el cielo nocturno y a cambiar su alumbrado.}{La astronomía debería enseñarse de noche.}{Los comerciantes se oponen a todo cambio.}
\sol{B. Cubre los cinco párrafos. (1.2)}
\item ¿Qué función cumple el cuarto párrafo?
\opcionesl{Presentar el problema inicial.}{Narrar el éxito de la primera noche.}{Mostrar la discusión que provocó el éxito.}{Concluir con la opinión de la profesora.}
\sol{C. (1.2)}
\item Según el texto, la alcaldía aceptó la propuesta
\opcionesl{con entusiasmo.}{a regañadientes y con vigilancia policial.}{solo para los niños.}{después de cambiar las lámparas.}
\sol{B. \S\,2. (1.1)}
\item Todas las siguientes reacciones de los asistentes se mencionan EXCEPTO
\opcionesl{ver por primera vez la Vía Láctea.}{llorar sin saber por qué.}{recordar historias de los padres.}{tomar fotografías con telescopios.}
\sol{D. (1.3)}
\item Los ingenieros argumentaron que las lámparas dirigidas hacia el suelo
\opcionesl{son más baratas de comprar.}{iluminan mejor la calle y gastan menos.}{reducen la delincuencia a la mitad.}{permiten ver la Vía Láctea todas las noches.}
\sol{B. \S\,4. (1.1)}
\item La expresión «las estrellas se habían vuelto un rumor de los abuelos» \pf{1} significa que
\opcionesl{los abuelos hablaban mal de las estrellas.}{las estrellas casi no se veían y solo los mayores las recordaban.}{los abuelos sabían astronomía.}{las estrellas hacían ruido.}
\sol{B. (1.1)}
\end{ejercicios}
